\documentclass[11pt,a4paper]{article}

% =============================================================
% ANONYMIZATION TOGGLE
% Set \anonymizedtrue  for double-blind submission
% Set \anonymizedfalse for arXiv, preprint, camera-ready
% Change ONLY THIS LINE:
\newif\ifanonymized
\anonymizedfalse
% =============================================================

% --- Core packages ---
\usepackage[utf8]{inputenc}
\usepackage[T1]{fontenc}
\usepackage{lmodern}
\usepackage[english]{babel}
\usepackage{geometry}
\geometry{margin=1in}
\usepackage{graphicx}
\usepackage{booktabs}
\usepackage{amsmath}
\usepackage{cleveref}
\usepackage{enumitem}
\usepackage{xcolor}
\usepackage{listings}
\usepackage{authblk}
\usepackage{microtype}
\usepackage{tabularx}
\usepackage{longtable}

\usepackage{hyperref}
\ifanonymized
  \hypersetup{
    colorlinks=true,
    linkcolor=blue!70!black,
    citecolor=blue!70!black,
    urlcolor=blue!70!black,
    pdfauthor={},
    pdftitle={},
  }
\else
  \hypersetup{
    colorlinks=true,
    linkcolor=blue!70!black,
    citecolor=blue!70!black,
    urlcolor=blue!70!black,
  }
\fi

\lstset{
  basicstyle=\small\ttfamily,
  breaklines=true,
  frame=single,
  columns=fullflexible,
  aboveskip=8pt,
  belowskip=8pt,
}

% =============================================================
% TITLE AND AUTHORS
% Replace the placeholders below with your paper's content.
% =============================================================

\title{[Paper Title: Concise, Specific, Informative]}

\ifanonymized
  \author{[Anonymous for review]}
  \date{}
\else
  \author{William Oliveira\\
  Independent Researcher\\
  ORCID: 0000-0000-0000-0000\\
  \texttt{contact@woliveiras.com}}
  \date{[Month YYYY]}
\fi

% =============================================================

\begin{document}

% --- Title page ---
\begin{titlepage}
  \centering
  \vspace*{\fill}
  {\LARGE\bfseries [Paper Title]\par}
  \vspace{2em}
  \ifanonymized
    {\large [Anonymous for review]\par}
  \else
    {\large William Oliveira\par}
    \vspace{0.5em}
    {\normalsize Independent Researcher\par}
    {\normalsize ORCID: 0000-0000-0000-0000\par}
    {\normalsize \texttt{contact@woliveiras.com}\par}
    \vspace{1em}
    {\normalsize [Month YYYY]\par}
  \fi
  \vspace*{\fill}
\end{titlepage}

% --- Abstract page ---
\begin{titlepage}
  \vspace*{\fill}
  \begin{abstract}
% Structured abstract: Background / Objective / Method / Results / Conclusion
% Background: What is the context? Why does this problem exist?
% Objective: What is the specific goal of this paper?
% Method: What study type and approach did you use? (≤2 sentences)
% Results: What did you find? (include at least ONE specific result)
% Conclusion: What does this mean? What should practitioners/researchers do?
[Background: context and motivation in 1–2 sentences.]
[Objective: the specific goal of this paper in 1 sentence.]
[Method: study design and approach in 1–2 sentences.]
[Results: key findings - include at least one specific result.]
[Conclusion: takeaway and implications in 1–2 sentences.]
  \end{abstract}
  \vspace*{\fill}
\end{titlepage}

% --- Table of contents (remove for conference papers) ---
\tableofcontents
\thispagestyle{empty}
\newpage

% =============================================================
\section{Introduction}
\label{sec:introduction}
% =============================================================
% Checklist:
% [ ] Problem is concrete (not "X is important")
% [ ] Gap is supported with evidence (cite or explain why gap exists)
% [ ] Research questions are numbered and scoped
% [ ] Contributions are numbered and match what is delivered
% [ ] Paper organization paragraph at the end

[Motivation: describe the problem and its practical impact. Be concrete.]

[Gap: what existing work does NOT address. Support with citations.]

We investigate the following research questions:

\begin{itemize}
  \item \textbf{RQ1:} [First research question - specific, measurable, scoped.]
  \item \textbf{RQ2:} [Second research question.]
  % \item \textbf{RQ3:} [Third research question, if applicable.]
\end{itemize}

Our contributions are:
\begin{enumerate}
  \item [First contribution - specific and verifiable.]
  \item [Second contribution.]
  % \item [Third contribution, if applicable.]
\end{enumerate}

The remainder of this paper is organized as follows.
\Cref{sec:background} reviews background and related work.
\Cref{sec:study-design} describes the study design.
\Cref{sec:results} presents the findings.
\Cref{sec:discussion} discusses implications.
\Cref{sec:threats} addresses threats to validity.
\Cref{sec:conclusion} concludes.


% =============================================================
\section{Background and Related Work}
\label{sec:background}
% =============================================================
% Checklist:
% [ ] Background covers concepts needed to understand the paper
% [ ] Related work describes what others did AND what they did NOT do
% [ ] Gap statement at end of related work section

\subsection{[Core Concept 1]}
\label{sec:background-concept1}

[Background explanation. Cite foundational sources~\cite{FIXME}.]

\subsection{[Core Concept 2]}
\label{sec:background-concept2}

[Background explanation.]

\subsection{Related Work}
\label{sec:related-work}

[Summary of related work - organized by theme or approach. For each group of papers, state what they study and what they do NOT address.]

% Positioning table (uncomment and fill for ≥5 related papers):
% \begin{table}[ht]
% \centering
% \small
% \caption{Comparison of related work. ✓ = studied, △ = partially, ✗ = not studied.}
% \label{tab:related-work}
% \begin{tabular}{lccc}
% \toprule
% Study & Dimension A & Dimension B & Dimension C \\
% \midrule
% Paper A~\cite{A} & ✓ & △ & ✗ \\
% Paper B~\cite{B} & ✓ & ✓ & ✗ \\
% \textbf{This paper} & ✓ & ✓ & ✓ \\
% \bottomrule
% \end{tabular}
% \end{table}

Unlike prior work, this paper addresses [gap] by [approach].


% =============================================================
\section{Study Design}
\label{sec:study-design}
% =============================================================
% Checklist:
% [ ] Study type is named (case study / benchmarking / experiment / SLR / survey)
% [ ] Subjects/objects described with enough detail for replication judgment
% [ ] Data collection method and instruments described
% [ ] Analysis procedure described (how categories/metrics were derived)

\subsection{Study Type and Goals}

This paper presents [a benchmarking study / a practitioner case study / a controlled experiment / ...].
The goal is to [GQM: analyze X for the purpose of Y with respect to Z from the viewpoint of W in the context of V].

\subsection{Subjects / Experimental Units}

[Describe what was studied: hardware, models, system, dataset, participants - whatever is the unit of study. Include exact versions, sizes, specifications.]

\subsection{Data Collection}

[Describe how data was collected: instruments, measurements, sources. For benchmarking: specify N per condition, warm-up runs, run order, thermal policy, system isolation. For case study: data sources, classification criteria.]

\subsection{Analysis Procedure}

[Describe how raw data was analyzed: statistical tests (justify choices), coding scheme for qualitative data, aggregation method.]

\subsection{Replication Package}

\ifanonymized
  The replication package (data, scripts, and analysis notebooks) will be made available upon acceptance at [anonymous repository link].
\else
  The replication package is available at: \url{https://github.com/[FIXME]/[FIXME]}.
\fi


% =============================================================
\section{Results}
\label{sec:results}
% =============================================================
% Checklist:
% [ ] Each subsection answers one RQ
% [ ] Subsection title states the finding, not just the question
% [ ] Figures/tables are referenced in text before they appear
% [ ] Mean AND variance reported for all quantitative results
% [ ] Statistical tests and p-values + effect sizes reported

\subsection{RQ1: [Restate the question]}
\label{sec:rq1}

% Tip: The subsection title should encode the answer, e.g.:
% "RQ1: On Radxa Rock 5B, throughput scales sublinearly with model size"

[Present findings for RQ1. Reference figures and tables.]

% Example table:
% \begin{table}[ht]
% \centering
% \small
% \caption{[Description of what the table shows].}
% \label{tab:rq1-results}
% \begin{tabularx}{\textwidth}{lXX}
% \toprule
% \textbf{Condition} & \textbf{Mean ± SD} & \textbf{[Metric]} \\
% \midrule
% Condition A & 45.2 ± 3.1 & ... \\
% Condition B & 38.7 ± 2.8 & ... \\
% \bottomrule
% \end{tabularx}
% \end{table}

\textbf{Answer to RQ1:} [1–2 sentence direct answer to the research question.]

\subsection{RQ2: [Restate the question]}
\label{sec:rq2}

[Present findings for RQ2.]

\textbf{Answer to RQ2:} [1–2 sentence direct answer.]


% =============================================================
\section{Discussion}
\label{sec:discussion}
% =============================================================
% Checklist:
% [ ] Interpretation of results (why do findings look like this?)
% [ ] Comparison to related work (agrees / contradicts / extends?)
% [ ] Implications for practitioners (actionable recommendations)
% [ ] Implications for researchers (open questions, future work directions)
% [ ] Limitations (honest scope bounds - different from threats to validity)

\subsection{Interpretation}

[Why do the results look like this? Provide explanations grounded in the data.]

\subsection{Comparison to Related Work}

[How do findings agree or disagree with related work? If they disagree, explain why.]

\subsection{Implications for Practitioners}

\begin{enumerate}
  \item \textbf{[Heuristic / recommendation 1]:} [Explanation.]
  \item \textbf{[Heuristic / recommendation 2]:} [Explanation.]
\end{enumerate}

\subsection{Implications for Researchers}

[What questions does this work open? What follow-on studies are needed?]


% =============================================================
\section{Threats to Validity}
\label{sec:threats}
% =============================================================
% Use Wohlin framework: Construct / Internal / External / Conclusion
% For each threat: name it → describe it → state your mitigation (or acknowledge if unmitigated)

\subsection{Construct Validity}

[Threats to whether measured constructs match intended theoretical constructs.]

\subsection{Internal Validity}

[Threats to whether observed effects are caused by the treatment (experiments, case studies).]

\subsection{External Validity}

[Threats to generalizability. Be specific about context factors that bound applicability.]

\subsection{Conclusion Validity}

[Threats to whether statistical conclusions are correct: power, test assumptions, multiple comparisons.]


% =============================================================
\section{Conclusion}
\label{sec:conclusion}
% =============================================================
% Checklist:
% [ ] Each RQ answered explicitly (1 sentence per RQ)
% [ ] No new claims not present in the paper body
% [ ] Future work is specific (not "more research is needed")

[Summary of motivation and approach in 2–3 sentences.]

In answer to our research questions:
\textbf{RQ1:} [Direct answer.]
\textbf{RQ2:} [Direct answer.]

[Broader takeaway: what does this mean for the field?]

Future work includes [specific next step 1] and [specific next step 2].


% =============================================================
% BIBLIOGRAPHY
% =============================================================

\bibliographystyle{plain}
\bibliography{references}
% Or with biblatex: \printbibliography


% =============================================================
% DATA AVAILABILITY (required by most venues)
% =============================================================
\section*{Data Availability}
\ifanonymized
  Data and scripts will be made publicly available upon acceptance.
\else
  Data, scripts, and analysis notebooks are available at: \url{https://github.com/[FIXME]/[FIXME]}
  under the MIT license (scripts) and CC BY 4.0 (paper).
\fi


% =============================================================
% ACKNOWLEDGMENTS (hidden in blind submissions)
% =============================================================
\ifanonymized
  % [Acknowledgments hidden for blind review]
\else
  \section*{Acknowledgments}
  [Funding sources, reviewers, etc. If none: This research received no external funding.]
\fi

\end{document}
